\chapter{Аналитический раздел}

\section{Предметная область}

Совместное редактирование данных — это подход к организации работы с информацией, при котором несколько пользователей имеют возможность одновременно просматривать, изменять и дополнять общее содержимое. Для обеспечения согласованности данных в таких системах применяются специальные протоколы взаимодействия, которые регламентируют порядок обмена изменениями, их синхронизацию и разрешение конфликтов между параллельными правками.

Протокол совместного редактирования заметок представляет собой набор правил и механизмов, обеспечивающих передачу, обработку и согласование изменений текстовых данных между клиентами и сервером (или между узлами распределённой системы). Основной задачей такого протокола является поддержание целостности и актуальности заметок при одновременной работе нескольких пользователей, а также минимизация потерь данных и конфликтных ситуаций.

В рамках применения протокола совместного редактирования заметок предусматривается обмен следующими типами данных: идентификаторы заметок и пользователей, содержимое изменений, версии документа, временные метки операций и служебная информация, необходимая для синхронизации. Протокол должен обеспечивать корректное применение изменений независимо от порядка их поступления и учитывать возможные задержки передачи данных по сети.

Использование протокола совместного редактирования особенно актуально в следующих ситуациях:

\begin{itemize}
    \item \textbf{совместная работа над текстовыми материалами:}
    при командной разработке документов, заметок, технической или учебной документации, когда несколько пользователей одновременно вносят изменения и требуется видеть актуальное состояние текста без необходимости ручного объединения правок;
    \item \textbf{удалённое взаимодействие пользователей:}
    при работе распределённых команд, находящихся в разных географических точках, протокол совместного редактирования позволяет обеспечить синхронную работу с заметками в режиме реального времени через сеть Интернет;
    \item \textbf{обучающие и исследовательские процессы:}
    в учебных и научных проектах протокол совместного редактирования может использоваться для коллективного ведения конспектов, заметок и обсуждений, повышая эффективность взаимодействия участников;
    \item \textbf{личные и корпоративные системы управления знаниями:}
    в системах заметок, таск-трекерах и внутренних базах знаний протокол совместного редактирования позволяет нескольким пользователям актуализировать информацию без риска её утраты или рассинхронизации.
\end{itemize}

Протокол совместного редактирования заметок обеспечивает согласованность данных, возможность оперативного обмена изменениями и устойчивость системы к сетевым задержкам и конфликтам. За счёт этого достигается повышение удобства работы пользователей, сокращение времени на согласование правок и повышение надёжности хранения информации.

В целом, применение протоколов совместного редактирования расширяет возможности традиционных систем заметок, позволяя организовать эффективную коллективную работу с текстовой информацией в реальном времени, в том числе в условиях распределённых и асинхронных взаимодействий.

\section{Возможные форматы применения протокола и целевая аудитория}

Разработка протокола для совместного редактирования заметок ориентирована на широкий спектр сфер, в которых требуется коллективная работа с текстовой информацией. В рамках данной курсовой работы рассматриваются основные форматы применения разработанного протокола, а также ключевые группы целевой аудитории, которые могут извлечь практическую пользу из его использования.

Во всех рассматриваемых форматах протокол совместного редактирования позволяет повысить эффективность взаимодействия пользователей, обеспечить согласованность данных, сократить время на обмен информацией и минимизировать ошибки, связанные с параллельным внесением изменений.

\subsection*{Совместная работа над учебными и учебно-методическими материалами}


Протокол совместного редактирования заметок может применяться в образовательной среде для коллективного ведения конспектов, методических материалов и учебных документов. Несколько пользователей могут одновременно работать с одной заметкой, вносить правки, дополнять информацию и видеть актуальное состояние документа в режиме реального времени.

Основными пользователями данного формата являются следующие категории:

\begin{itemize}
    \item студенты, совместно работающие над конспектами лекций и практических занятий;
    \item преподаватели, разрабатывающие учебно-методические материалы;
    \item учебные группы и проектные команды в образовательных учреждениях;
    \item участники онлайн-курсов и дистанционного обучения.
\end{itemize}

\subsection*{Командная работа в распределённых и удалённых командах}

В условиях удалённой работы и распределённых команд протокол совместного редактирования заметок позволяет организовать эффективное взаимодействие между участниками проекта, находящимися в разных географических точках. Заметки и рабочие документы синхронизируются автоматически, что исключает необходимость ручного объединения версий и пересылки файлов.

Основными пользователями данного формата являются следующие категории:

\begin{itemize}
    \item IT-команды и разработчики программного обеспечения;
    \item проектные и продуктовые менеджеры;
    \item распределённые рабочие группы и стартап-команды;
    \item сотрудники компаний, использующие удалённый или гибридный формат работы.
\end{itemize}

\subsection*{Ведение личных и корпоративных баз знаний}

Протокол совместного редактирования заметок может использоваться в системах управления знаниями для коллективного накопления и актуализации информации. Это особенно важно в корпоративной среде, где несколько сотрудников одновременно работают с внутренними регламентами, инструкциями и справочными материалами.

Основными пользователями данного формата являются следующие категории:

\begin{itemize}
    \item корпоративные пользователи и сотрудники организаций;
    \item команды технической поддержки и сопровождения;
    \item специалисты по обучению и адаптации персонала;
    \item администраторы внутренних баз знаний и документации.
\end{itemize}

\subsection*{Совместное планирование и управление проектами}

В системах планирования и управления проектами протокол совместного редактирования заметок позволяет участникам проекта совместно вести проектную документацию, списки задач, протоколы встреч и рабочие заметки. Это обеспечивает актуальность информации и упрощает контроль за ходом выполнения проекта.

Основными пользователями данного формата являются следующие категории:

\begin{itemize}
    \item руководители проектов;
    \item члены проектных команд;
    \item бизнес-аналитики;
    \item консультанты и внешние подрядчики.
\end{itemize}

\subsection*{Научно-исследовательская и аналитическая деятельность}

В научной и исследовательской деятельности протокол совместного редактирования заметок может применяться для коллективной работы над аналитическими материалами, гипотезами, результатами экспериментов и промежуточными выводами. Совместное редактирование снижает вероятность расхождений в данных и повышает прозрачность работы исследовательской группы.

Основными пользователями данного формата являются следующие категории:

\begin{itemize}
    \item научные сотрудники и исследователи;
    \item аспиранты и магистранты;
    \item участники научных и исследовательских проектов;
    \item аналитические группы и исследовательские центры.
\end{itemize}


\section{Протоколы передачи данных для совместного редактирования заметок}

\subsection*{WebSocket}

WebSocket — это протокол, обеспечивающий двусторонний постоянный канал связи между клиентом и сервером поверх TCP. Он широко используется в веб-приложениях для передачи данных в режиме реального времени.

Преимущества:

\begin{enumerate}
    \item поддержка постоянного соединения и двустороннего обмена данными;
    \item низкие задержки при передаче изменений;
    \item эффективное использование сетевых ресурсов по сравнению с частыми HTTP-запросами;
    \item хорошая поддержка в современных браузерах и серверных платформах;
    \item удобен для реализации синхронного совместного редактирования.
\end{enumerate}


\subsection*{HTTP(S) (Hypertext Transfer Protocol (Secure))}

HTTP — широко распространённый протокол передачи данных в сети Интернет. При использовании HTTPS обеспечивается шифрование передаваемой информации.

Преимущества:

\begin{itemize}
    \item универсальность и совместимость с существующей веб-инфраструктурой;
    \item простота реализации и отладки;
    \item широкая поддержка библиотек и инструментов;
    \item высокая безопасность при использовании HTTPS.
\end{itemize}

Недостатки:

\begin{itemize}
    \item не поддерживает постоянное двустороннее соединение по умолчанию;
    \item требует частых запросов для получения обновлений;
    \item менее эффективен для задач реального времени по сравнению с WebSocket.
\end{itemize}

\subsection*{WebSocket поверх HTTP(S)}

Данный подход предполагает установку WebSocket-соединения через HTTP(S)-рукопожатие, что позволяет совместить преимущества обоих протоколов.

Преимущества:

\begin{itemize}
    \item совместимость с существующей веб-инфраструктурой;
    \item безопасная передача данных при использовании TLS;
    \item высокая скорость синхронизации изменений;
    \item удобство интеграции в веб-приложения.
\end{itemize}

Недостатки:

\begin{itemize}
    \item усложнение архитектуры приложения;
    \item повышенные требования к серверным ресурсам.
\end{itemize}

\subsection*{gRPC}

gRPC — это высокопроизводительный протокол удалённого вызова процедур, основанный на HTTP/2 и бинарной сериализации данных.

Преимущества:

\begin{itemize}
    \item высокая производительность и низкие задержки;
    \item поддержка потоковой передачи данных;
    \item строгая типизация сообщений;
    \item эффективен для взаимодействия между микросервисами.
\end{itemize}

Недостатки:

\begin{itemize}
    \item сложнее в реализации по сравнению с HTTP и WebSocket;
    \item ограниченная поддержка в браузерах без дополнительных прокси;
    \item избыточен для простых систем заметок.
\end{itemize}

\section{Выбор протокола}

\subsection*{Выбор транспортного протокола}

При разработке протокола уровня приложений для совместного редактирования текстовых заметок в качестве транспортного уровня рассматриваются протоколы TCP (Transmission Control Protocol) и UDP (User Datagram Protocol). Оба протокола широко применяются для передачи данных в компьютерных сетях, однако обладают различными характеристиками, которые имеют принципиальное значение в контексте разрабатываемого протокола.

TCP является одним из базовых транспортных протоколов сети Интернет и обеспечивает надёжную, упорядоченную и гарантированную доставку данных между приложениями, работающими на различных узлах сети. TCP функционирует совместно с протоколом IP (Internet Protocol) и в совокупности с ним составляет основу большинства сетевых приложений.

Основные свойства TCP:

\begin{itemize}
    \item гарантированная доставка данных между клиентом и сервером;
    \item контроль потока и повторная передача данных при их потере;
    \item упорядоченность сообщений и отсутствие дублирования пакетов;
    \item более высокая нагрузка на сервер за счёт необходимости поддержания соединений;
    \item подходит для передачи критически важных данных, где требуется целостность и согласованность информации.
\end{itemize}

UDP, в отличие от TCP, ориентирован на простоту и минимальные задержки передачи данных, а не на надёжность и упорядоченность. Он также работает поверх IP, однако предоставляет значительно более базовый уровень сервиса и не гарантирует доставку пакетов.

Основные свойства UDP:

\begin{itemize}
    \item минимальная задержка передачи данных;
    \item отсутствие гарантии доставки и упорядоченности сообщений;
    \item отсутствие механизмов повторной передачи данных;
    \item подходит для приложений реального времени, где допустимы небольшие потери информации;
    \item снижает нагрузку на сервер за счёт отсутствия необходимости установления соединения.
\end{itemize}

В рамках задачи совместного редактирования текстовых заметок ключевыми требованиями являются целостность данных, корректная синхронизация версий документа и отсутствие потерь или дублирования операций редактирования. Потеря или нарушение порядка передачи сообщений может привести к рассинхронизации состояния документа и возникновению конфликтов между версиями текста.

Исходя из требований к надёжности и согласованности данных, в качестве транспортного протокола был выбран TCP. Использование TCP позволяет гарантировать доставку сообщений, их упорядоченность и целостность, что является критически важным для корректной работы протокола совместного редактирования и обеспечения согласованного состояния текстовых заметок у всех участников процесса.

\subsection*{Выбор алгоритма шифрования}

При разработке протокола уровня приложений для совместного редактирования текстовых заметок необходимо обеспечить защиту передаваемых данных от несанкционированного доступа и модификации. В рамках данной задачи рассматриваются наиболее распространённые криптографические алгоритмы, применяемые для защиты сетевых взаимодействий, а именно AES (Advanced Encryption Standard), ChaCha20 и RSA (Rivest–Shamir–Adleman).

Передаваемые в рамках протокола данные включают аутентификационную информацию пользователей, содержимое текстовых заметок, а также служебные сообщения, связанные с синхронизацией версий и передачей операций редактирования. Нарушение конфиденциальности или целостности этих данных может привести к утечке информации или искажению состояния документов.

AES является симметричным блочным алгоритмом шифрования, широко используемым для защиты конфиденциальных данных. В настоящее время AES является международным стандартом де-факто и де-юре и применяется в большинстве современных протоколов безопасности.

Преимущества AES:

\begin{itemize}
    \item высокий уровень криптографической стойкости к современным атакам;
    \item высокая производительность при шифровании и дешифровании данных;
    \item высокая производительность при шифровании и дешифровании данных;
    \item широкая поддержка в криптографических библиотеках и языках программирования;
    \item хорошо подходит для шифрования потоков данных в клиент-серверных приложениях.
\end{itemize}

Недостатки AES:

\begin{itemize}
    \item требует безопасного обмена симметричным ключом между участниками взаимодействия.
\end{itemize}

ChaCha20 является симметричным поточным алгоритмом шифрования, разработанным как альтернатива AES, особенно в средах, где отсутствует аппаратная поддержка AES-NI.

Преимущества ChaCha20:

\begin{itemize}
    \item высокая скорость работы на процессорах без аппаратного ускорения;
    \item устойчивость к атакам по сторонним каналам;
    \item относительная простота реализации.
\end{itemize}

Недостатки ChaCha20:

\begin{itemize}
    \item менее широко используется по сравнению с AES;
    \item в средах с аппаратной поддержкой AES может уступать ему по производительности.
\end{itemize}

RSA является асимметричным алгоритмом шифрования, основанным на сложности факторизации больших чисел. Он широко используется для безопасного обмена ключами и цифровых подписей, но не предназначен для шифрования больших объёмов данных.

Преимущества RSA:

\begin{itemize}
    \item возможность безопасного обмена ключами без предварительного защищённого канала;
    \item используется для аутентификации и цифровых подписей;
    \item хорошо подходит для начального этапа установления защищённого соединения.
\end{itemize}

Недостатки RSA:

\begin{itemize}
    \item значительно более низкая производительность по сравнению с симметричными алгоритмами;
    \item неэффективен для шифрования больших объёмов данных;
    \item требует использования длинных ключей (2048 бит и более) для обеспечения высокой безопасности.
\end{itemize}

С учётом требований к безопасности, производительности и простоте интеграции в клиент-серверную архитектуру протокола совместного редактирования текстовых заметок, в качестве основного алгоритма шифрования целесообразно использовать AES для защиты передаваемых данных. Асимметричный алгоритм RSA может применяться на этапе установления соединения или обмена симметричными ключами, после чего дальнейшая передача данных осуществляется с использованием AES.

Выбор AES позволяет обеспечить высокий уровень защиты информации, минимальные накладные расходы при передаче данных и эффективную работу протокола совместного редактирования в условиях одновременной работы нескольких пользователей.

\section{Постановка задачи}

На основе проведённого анализа предметной области, рассмотренных форматов применения протокола совместного редактирования заметок, а также обзора существующих протоколов передачи данных, перед нами стоит задача разработки протокола уровня приложений для совместного редактирования текстовых заметок.

Разрабатываемый протокол должен обеспечивать корректное и согласованное взаимодействие участников процесса совместного редактирования, поддерживать синхронизацию версий документа, передачу операций редактирования и разрешение конфликтов, возникающих при одновременном изменении текста несколькими пользователями.

В рамках данной курсовой работы протокол должен удовлетворять следующим требованиям:

\begin{itemize}
    \item протокол должен обеспечивать синхронизацию версий текстовой заметки между всеми участниками редактирования, исключая рассинхронизацию и потерю данных;
    \item протокол должен поддерживать передачу операций редактирования (добавление, удаление и изменение текста), а также их корректное применение независимо от порядка поступления сообщений;
    \item протокол должен предусматривать механизмы разрешения конфликтов при одновременном редактировании текста несколькими пользователями;
    \item протокол должен предусматривать аутентификацию участников процесса совместного редактирования и защиту передаваемых данных от несанкционированного доступа и модификации;
\end{itemize}

Исходя из перечисленных требований, целью данной курсовой работы является разработка протокола уровня приложений для совместного редактирования текстовых заметок, а также реализация и тестирование программного приложения, использующего разработанный протокол. Разработанный протокол должен обеспечивать согласованность данных, корректную синхронизацию изменений и удобство коллективной работы пользователей в условиях распределённой среды.
